\section{Stage 2: Multi-Source Candidate Ranking}
\label{sec:lit-multi-source}

As mentioned earlier, most commercial ATS platforms and the majority of research literature still rely
on a self-reported CV as the only source of evidence for a candidate. Research on resume fraud and embellishment
is well documented~\cite{henle2019_resume_fraud,dunning2004}, meanwhile employer surveys report that qualified
applicants are screened out when they do not match the job-description wording verbatim~\cite{fuller2021hiddenworkers}.
While approaches such as keyword-matching and semantic analysis can rank candidates, it never truly
questions whether the candidate actually has the skills they claim to have.

\subsection{Role Classification from Developer Platforms}
Montandon et al.~\cite{montandon2021_roles} explore the use of GitHub as a supplementary source of evidence,
rather than only the CV. They leverage GitHub profiles and repository metadata in order to infer a 
total of six different technical roles. They use programming language usage, descriptions and bios
to train a supervised machine learning model to classify the role of a candidate. The result of this is 
an aggregate precision of 0.88 and an area under the curve (AUC) of 0.89. The labelled set is comprised
of 2,284 developers, with language features being the most discriminative signal.
For most roles, especially DevOps and mobile, the recall is weaker. Authors treat the classifier as a 
supplementary hiring signal rather than a sole criterion understandably, since this would be a very rigid 
approach to screening.

The approach still remains limited. Profile text is matched with regex-style cues, vulnerable to 
self-report bias, also, public activity may not reflect candidates who work in private repositories.
Alongside this, role labels are fixed and coarse, rather than JD-configurable, and finally predictions
are not cross-checked against the CV, and neither are they explained in auditable terms.
Regardless of this, the study motivated the treatment of GitHub as a supplementary source of evidence,
rather than just being an isolated label among a fixed set of roles.

\subsection{Repository Mining and Enrichment}
Recent work pushes further away from the CV-only dependency, exploring the use of GitHub as a supplementary 
source of evidence. First, we review Vaishampayan et al.~\cite{vaishampayan2025_gitmeter}, who develop GitMeter.
Their system mines public repositories to discover essential candidate data such as technology use and code quality.
Their user study of 20 stakeholders reports that they are willing to use repository evidence in shortlisting,
highlighting the potential of repository evidence in hiring decisions. Later, Kumar et al.~\cite{kumar2025_evidence_hiring}
extend this line of work with broad language analysis (through CodeQL), plagiarism and AI-detection tooling.
They also enrich the data with academic profiles, providing a more comprehensive view of a candidate's skills and experience.

\label{sec:contrastive-github-hiring}
While their work moves beyond the CV-only limitation, GitMeter~\cite{vaishampayan2025_gitmeter} augments resume matching with GitHub
evidence rather than replacing the CV altogether. Kumar et al.~\cite{kumar2025_evidence_hiring} push further towards GitHub-centric
assessment enriched with academic profiles. Neither approach fuses CV, GitHub, and LinkedIn per metric with traceable conflict resolution,
nor do they report recruiter-configurable, job-description-driven metric weights or per-source attributions of an overall score.
This work was identified after the specification and implementation of MERIT, however, we still include it for comparison and to
highlight the gap that MERIT is designed to solve. Note that we did not use this work to derive the requirements, architecture or
evaluation in this dissertation.

\subsection{The Requirement for Evidence Verification}
The gap that motivates MERIT's multi-source pillar is therefore not merely ``add GitHub,'' but \emph{verify} 
CV claims with heterogeneous evidence, quantify uncertainty when sources disagree, and rank candidates on a common scale.

\section{Bayesian Evidence Fusion}

A key challenge in developing MERIT will be to address the uncertaity in conflicting signals that arise from handling multi-source candidate data.
MERIT employs a probabilistic framework based on Bayesian Evidence Fusion to address this challenge.
This is in contrast to traditional-rule based systems that rely on single sources of information.
By leveraging multiple sources, MERIT is able to provide a thorough "Glass-Box" assessment to separate the estimated ability from the certainty of that estimate.

\subsection*{The Beta Distribution Model}
We use a Beta distribution denoted as $Beta(\alpha, \beta)$ to model the "probability of competence" for a specific skill.
This distribution is well-suited for the task since its distribution can be updated incrementally as new evidence arrives.
This is key to providing an extensible, modular architecture that allows for the integration of new evidence sources beyond those that will be implemented in this project.

When we calculate the beta distribution for a metric, we can expect the following:
\begin{enumerate}
    \item \textbf{The Match Score:} The expected value (mean) of the distribution, providing the most likely estimate of a candidate's skill level.
    \item \textbf{The Confidence Level:} The spread (variance) of the distribution; a narrow distribution (low variance) signifies high confidence in the data, while a wide distribution (high variance) indicates a lack of evidence or conflicting signals.
\end{enumerate}

The parameters of the distribution are updated based on the type of evidence discovered:
\begin{itemize}
    \item \textbf{Alpha ($\alpha$):} Represents "success" signals or positive evidence (e.g., a verified GitHub contribution, a high-impact project on a CV).
    \item \textbf{Beta ($\beta$):} Represents "failure" signals, noise, or uncertainty (e.g., skill decay over time, lack of evidence for a self-reported claim).
\end{itemize}

The fused score (Match Score) is calculated as the mean of the distribution \cite{gelman2013bayesian}:
\begin{equation}
    E[X] = \frac{\alpha}{\alpha + \beta}
\end{equation}

The confidence level is calculated as the variance of the distribution:
\begin{equation}
    Var[X] = \frac{\alpha \beta}{(\alpha + \beta)^2 (\alpha + \beta + 1)}
\end{equation}

\subsection*{Modelling Uncertainty}
A key advantage of this framework is the ability to quantify \textit{uncertainty} through the variance of the distribution. 
As more evidence is gathered, the distribution narrows, increasing the confidence in the assessment. 
This allows MERIT to flag candidates with "High Confidence" vs. "Low Confidence" labels, providing the recruiter with a measure of how much to trust the automated ranking. 
By utilising the conjugate prior properties of the Beta distribution, MERIT can efficiently update its beliefs as new data sources are ingested.

